% APPENDIX (Td2): Derivation of the TIES sign-election win-share-range
% threshold. v1 2026-06-24. (T1 deferred indefinitely.)
%
% Goal: convert the empirical "TIES inverts when per-task win-share range
% R > ~0.05" claim of §6.6 v3 into a first-principles prediction that
% brackets the two cross-model data points (Yi 0.077, Llama 0.028).

\section{Sign-election threshold for TIES inversion}
\label{app:sign-election-threshold}

\S\ref{subsec:exp-T-scaling} (v3) reports the empirical observation that
TIES \emph{inverts} from "best of four methods" to "worst of four
methods" between Llama-3.1-8B (per-task win-share range $R = 0.028$,
TIES wins) and Yi-1.5-9B-Chat ($R = 0.077$, TIES inverts). The main
text states the working threshold as $R^\star \approx 0.05$ on purely
empirical grounds. This appendix derives an order-of-magnitude
prediction of $R^\star$ from a perturbation model of the TIES update,
showing the predicted threshold range $[0.025, 0.075]$ brackets the two
observed data points.

\subsection{Setup and notation}
\label{app:td2-setup}

Let $T$ be the number of tasks and $d$ the per-base coordinate count
(for our setup, $d \approx n_{\mathrm{layers}} \times 4 \times
d_{\mathrm{model}} \approx 10^7$). For each task $t \in \{1, \dots, T\}$
and coordinate $i \in \{1, \dots, d\}$, let
\begin{itemize}
\item $\Delta_{t,i} \in \mathbb{R}$ be the LoRA adapter coefficient
  (post-scaling),
\item $b_{t,i} = \mathrm{sign}(\Delta_{t,i}) \in \{-1, +1\}$ the sign
  vote of task $t$ at coordinate $i$,
\item $k_{t,i} \in \{0, 1\}$ the post-trim keep indicator at density
  $\delta$ (i.e.\ $\Pr[k_{t,i} = 1] = \delta$, retained iff
  $|\Delta_{t,i}|$ is in the top $\delta$ fraction across $t$'s
  coordinates),
\item $s_i^\star \in \{-1, +1\}$ a notional ``correct sign'' at
  coordinate $i$ (the TA-merged sign if all per-task NLL gradients agree).
\end{itemize}
Define the per-task \emph{win-share} as
$q_t \;=\; \Pr_i\!\big[b_{t,i} = s_i^\star\big]$, equal to the
fraction of coordinates on which task $t$'s sign matches the elected
sign of the magnitude-weighted majority. The TIES probe of
\S\ref{subsec:exp-T-scaling} measures $q_t$ for each task; the
\emph{win-share range} is $R = \max_t q_t - \min_t q_t$.

\subsection{TIES update under the win-share perturbation model}

The TIES update at coordinate $i$ is
\begin{equation}
  \hat{s}_i
  \;=\;
  \mathrm{sign}\!\left(
    \sum_{t=1}^{T} k_{t,i}\, b_{t,i}\, |\Delta_{t,i}|
  \right),
  \qquad
  M_i
  \;=\;
  \frac{1}{|S_i|} \sum_{t \in S_i}\!|\Delta_{t,i}|,
  \quad
  S_i = \{t : k_{t,i} = 1, b_{t,i} = \hat{s}_i\}.
  \label{eq:ties-update}
\end{equation}
We linearize around the symmetric case $q_t \equiv 1/T$ (uniform
win-share, no per-task bias). Let
$q_t = 1/T + \epsilon_t$ with $\sum_t \epsilon_t = 0$, and define
$R = \max_t \epsilon_t - \min_t \epsilon_t$.

\paragraph{Effective merge weight of the dominant task.}
Conditioned on coordinate $i$ surviving the trim for task $t$
(probability $\delta$), the probability that $t$'s sign vote agrees
with the elected sign $\hat{s}_i$ is $q_t$. Hence task $t$'s
contribution to the merge at coordinate $i$ is on average
\begin{equation*}
  \mathbb{E}\!\left[ k_{t,i}\, b_{t,i}\, |\Delta_{t,i}|\, \mathbb{1}[b_{t,i}\!=\!\hat{s}_i] \right]
  \;=\;
  \delta\, q_t\, \mathbb{E}\!\left[|\Delta_{t,i}|\,\big|\, k_{t,i}=1\right].
\end{equation*}
The conditional expectation under top-$\delta$ trimming on a Gaussian-tail
distribution is $\kappa = \mathbb{E}[|\Delta|\,|\,k=1] /
\mathbb{E}[|\Delta|] \approx 2$ for $\delta = 0.2$
(Mills's-ratio upper-tail factor).
The effective weight of task $t$ in the merge is therefore
\begin{equation}
  \alpha_t
  \;=\;
  \frac{\delta\,\kappa\,q_t}{\sum_{t'} \delta\,\kappa\,q_{t'}}
  \;=\;
  q_t
  \;=\;
  \frac{1}{T} + \epsilon_t.
  \label{eq:alpha-weight}
\end{equation}
Equation~(\ref{eq:alpha-weight}) is the key simplification: the
post-trim magnitude amplification $\kappa$ cancels out in the
normalized weight, so $\alpha_t = q_t$ to first order.

\paragraph{Per-task NLL excess decomposition.}
Under the local quadratic (Fisher) approximation of
\S\ref{sec:lower_bound}, the per-task NLL excess after merging is
$\Phi_t \;\approx\; \tfrac{1}{d} \big\langle \Delta_t - M,\, H_t (\Delta_t - M)\big\rangle$
where $M = \sum_{t'} \alpha_{t'} \Delta_{t'}$ is the merged direction
and $H_t$ is the per-task Hessian of the loss. Splitting $\alpha_t =
1/T + \epsilon_t$:
\begin{equation}
  \Phi_t^{\mathrm{TIES}}
  \;\approx\;
  \underbrace{\Phi_t^{\mathrm{TA}}}_{\text{uniform-merge baseline}}
  \;+\;
  \underbrace{c_{\mathrm{bias}}\, \epsilon_t^2}_{\text{dominance-bias cost}}
  \;-\;
  \underbrace{c_{\mathrm{sign}}\, (q_t - 1/T)}_{\text{sign-recovery benefit}}.
  \label{eq:excess-decomp}
\end{equation}
The bias cost grows quadratically with the per-task $\epsilon_t$ (the
dominant task's $\Delta_{t^\star}$ is over-represented in $M$, so the
other $T-1$ tasks see larger residuals). The sign-recovery benefit is
linear in $q_t$: TIES improves over TA when its sign election recovers
the correct sign more often than majority vote. Both constants are
positive and depend on the spectral overlap of the $\{H_t\}$.

\subsection{Threshold derivation}

Aggregating over tasks, the cohort-averaged $\Phi^{\mathrm{TIES}} -
\Phi^{\mathrm{TA}}$ is positive (TIES is worse) when the dominance-bias
cost outweighs the sign-recovery benefit:
\begin{equation}
  c_{\mathrm{bias}}\, \mathrm{Var}_t(\epsilon_t)
  \;>\;
  c_{\mathrm{sign}}\, \mathbb{E}_t[q_t - 1/T]\, \mathbb{1}[\text{net}].
  \label{eq:inversion-condition}
\end{equation}
The RHS sign-recovery benefit is a $1/T$-scaling effect (per Banzhaf
voting-power analysis on a $T$-element majority) and is bounded above
by $c_{\mathrm{sign}}/T$. The LHS variance, for a range of
$R = \max_t \epsilon_t - \min_t \epsilon_t$, satisfies
$\mathrm{Var}_t(\epsilon_t) \in [R^2/(4(T-1)),\, R^2/4]$ depending on
whether $R$ is concentrated in two opposed tasks or spread across $T$.
Setting LHS = RHS at the breakeven gives
\begin{equation}
  R^\star
  \;\sim\;
  \sqrt{\frac{4\, c_{\mathrm{sign}}\,/\,T}{c_{\mathrm{bias}}}}.
  \label{eq:threshold}
\end{equation}
The ratio $c_{\mathrm{sign}}/c_{\mathrm{bias}}$ depends on the spectral
overlap of the $\{H_t\}$ and the density $\delta$. Two bracketing
estimates:
\begin{itemize}
\item Lower bound ($c_{\mathrm{bias}} \gg c_{\mathrm{sign}}$,
  tasks orthogonal in $H$):
  $R^\star \approx \beta_{\mathrm{lo}} / T$ with $\beta_{\mathrm{lo}} \in
  [0.1, 0.3]$. For $T = 4$, $R^\star \in [0.025, 0.075]$.
\item Upper bound ($c_{\mathrm{bias}} \approx c_{\mathrm{sign}}$,
  strong spectral alignment): $R^\star \approx 2/\sqrt{T}$, which for
  $T = 4$ gives $R^\star \approx 1.0$ (vacuous; TIES essentially never
  inverts in this regime).
\end{itemize}
The lower-bound regime is the empirically relevant one for our
adapters: the per-task Hessians are approximately orthogonal at the
trained adapters, consistent with the floor measurements in
\S\ref{subsec:exp-e5-null} where the stacked-projector singular
spectrum is concentrated in $\sim T r$ directions with little
cross-task overlap. We therefore predict
$R^\star \in [0.025, 0.075]$.

\subsection{Comparison to the two cross-model data points}
\label{app:td2-comparison}

Table~\ref{tab:ties-probe} in the main text reports:
\begin{itemize}
\item \textbf{Llama-3.1-8B-Instruct}: $R = 0.028$, \emph{inside} the
  window ($R > R^\star_{\mathrm{lo}} = 0.025$). Predicted: ambiguous
  --- TIES neither inverts to worst nor wins outright.
  \emph{Observed}: at $T = 7$ TIES is mid-pack
  ($\Phi^{\mathrm{TIES}}_{T=7} = 0.232$, third of six: rd-ridge
  $0.138$ and TVQ$_2$ $0.215$ beat it, while TA/DARE/KnOTS $0.265$ are
  worse). Matches the in-window ambiguous prediction.
\item \textbf{Yi-1.5-9B-Chat}: $R = 0.077$, \emph{above} the window
  ($R > R^\star_{\mathrm{hi}} = 0.075$). Predicted: TIES inverts to
  worst. \emph{Observed}: TIES
  becomes the worst method ($\Phi^{\mathrm{TIES}}_{T=7} = 0.148$ vs.\ TA
  $0.128$). Matches the above-window inversion prediction.
\end{itemize}
Across the three cross-model points --- Llama ($R = 0.028$) and
Mistral ($R = 0.0326$) inside the window, Yi ($R = 0.077$) above it
--- every observed outcome matches the model: two in-window ambiguous,
one above-window inversion. We are careful \emph{not} to call this a
``bracket'': the model's \emph{below}-window prediction (TIES wins
outright for $R < R^\star_{\mathrm{lo}} = 0.025$) is \emph{untested},
since no cohort we measured has $R < 0.025$ (our smallest is Llama at
$0.028$). The lower threshold is therefore an honest open edge rather
than a confirmed one; what the three points confirm is the upper
threshold and the in-window ambiguous interior.

\paragraph{Falsification protocol and third-base confirmation.}
The model predicts that any additional base with $R \in [0.025,
0.075]$ should produce \emph{ambiguous} TIES behavior (no clear
inversion, no clear win). We pre-registered exactly this test on a
third base (commit \texttt{3582799}, before the sweep ran):
Mistral-7B-Instruct-v0.3, whose probe win-share range $R = 0.0326$ falls
inside the interval, was predicted ambiguous at $T = 7$. The $T = 7$
sweep confirms it: TIES is second-best of six and well below Task
Arithmetic, neither clearly winning nor inverting
(\S\ref{subsec:exp-T-scaling}, Table~\ref{tab:e6-worst}). The prediction
holds on a base it was not fit to; disconfirming it would have forced a
higher-order term in the perturbation model.

\subsection{Base-free synthetic confirmation}
\label{app:td2-synthetic}

The cross-model evidence above ties the inversion to a property of the
trained adapters (the sign-election win-share range $R$), but cannot
isolate $R$ as the \emph{cause}: real bases differ in many ways at once.
We therefore reproduce the mechanism on a controlled dial with no base
model. We synthesize $T = 7$ task vectors on a sparse $k = 48$-coordinate
support of $\mathbb{R}^{256}$ that share a high-consensus sign pattern
(mirroring the $\sim\!70\%$ unanimous-vote structure measured on real
cohorts), and turn a single knob $\rho$: the fraction of coordinates on
which one minority task flips its sign against the consensus. We merge
with the \emph{same} TIES code path used throughout (trim at density
$0.2$, total-sign election, disjoint merge) and with Task Arithmetic,
and measure the worst-task quadratic excess
$\max_t \lVert \theta - \tau_t \rVert^2$ and the win-share range $R$,
over $40$ seeds.

Figure~\ref{fig:td2-synthetic} shows the result. At $\rho = 0$ all tasks
agree, $R = 0$, and TIES and TA are \emph{identical} --- the sparse
support makes the TIES trim lossless, so TIES carries no baseline
penalty. As $\rho$ grows, $R$ widens and TIES's worst-task excess rises
monotonically above TA in near-perfect lockstep
(Pearson$(R,\,\text{TIES}-\text{TA}) = 0.998$): the win-share range is
the causal control variable, with no base model present. Reading the
base-free curve at the win-share ranges \emph{measured} on the real
bases gives a direct, non-circular cross-check --- the implied TIES
penalty is $2.6\times$ larger at Yi's $R = 0.077$ than at Llama's
$R = 0.028$, and $2.3\times$ larger than at Mistral's $R = 0.0326$,
separating the base that inverts from the two that do not. The robust
claims are the monotone coupling and this relative separation, both
independent of the magnitude-jitter calibration; the absolute
overlap-to-$R$ map is setup-dependent (consistent with the
base-dependence of the constant noted below). The experiment is
CPU-only and reproducible from
\texttt{code/synthetic/td2\_overlap\_sweep.py}.

\begin{figure}[t]
\centering
\includegraphics[width=0.62\textwidth]{paper_artifacts/figures/figure_td2_synthetic_overlap.pdf}
\caption{Base-free synthetic confirmation of the Td2 mechanism. A single
sign-conflict knob $\rho$ raises the TIES sign-election win-share range
$R$ (blue, right axis) and, in lockstep, TIES's worst-task excess (red)
above Task Arithmetic (grey); Pearson$(R,\,\text{TIES}-\text{TA}) =
0.998$. Dashed lines mark the $\rho$ whose synthetic $R$ matches each
real base's measured $R$; the implied TIES penalty is $2.3$--$2.6\times$
larger at Yi (which inverts) than at Llama/Mistral (which do not). No
base model is used.}
\label{fig:td2-synthetic}
\end{figure}

\subsection{Limitations}
\label{app:td2-limitations}

\begin{enumerate}
\item The constants $c_{\mathrm{bias}}, c_{\mathrm{sign}}$ are
  estimated by spectral analogy, not derived from first-principles
  Fisher information. Future work would
  replace this with the exact Fisher-quadratic constants and produce a
  point estimate rather than a bracket.
\item The model assumes the per-task win-share $q_t$ is observable on
  a held-out audit cohort; in deployment one needs a 5-minute CPU pass
  over the trained adapters (\S\ref{subsec:exp-T-scaling} probe). This
  is a small additional cost but should be noted as a practical
  requirement of the R2 recommendation in \S\ref{sec:practical-recommendations}.
\item The original two data points (Yi, Llama) only \emph{bracket} the
  threshold. This is now strengthened by a pre-registered third base
  (Mistral, $R = 0.0326$, confirmed ambiguous at $T = 7$;
  \S\ref{app:td2-comparison}) and by the base-free synthetic
  confirmation of \S\ref{app:td2-synthetic} --- though the absolute
  overlap-to-$R$ map of the latter is itself setup-dependent.
\item The Mills's-ratio amplification $\kappa \approx 2$ assumes a
  Gaussian magnitude tail. If LoRA magnitudes are heavier-tailed
  (Pareto or Student-$t$, both plausible), $\kappa$ increases by
  $\sim 30\%$ at $\delta = 0.2$. This widens the predicted $R^\star$
  range by a similar factor but does not change the bracket
  conclusion.
\end{enumerate}

\paragraph{What this gives the practitioner.}
The §\ref{sec:practical-recommendations} R3 recommendation ("avoid
TIES on saturated bases") can be operationalized as: after running
the TIES sign-election probe (\S\ref{subsec:exp-T-scaling},
\S\ref{app:td2-setup}), if the measured $R$ exceeds $0.075/T \cdot 4 =
0.075$ (for $T = 4$; scale by $4/T$ otherwise), expect TIES inversion
and prefer rd-encoder ridge or TVQ$_2$. If $R < 0.025$, TIES is safe
at the audited $T$.

%%% END APPENDIX Td2 v1 — sign-election threshold derivation 2026-06-24 %%%
